\documentclass[11pt]{article}
\usepackage{amsmath, amsthm, amssymb}
\usepackage[margin=1.1in]{geometry}
\usepackage{hyperref}

\newtheorem{theorem}{Theorem}
\newtheorem{proposition}[theorem]{Proposition}
\newtheorem{corollary}[theorem]{Corollary}
\newtheorem{lemma}[theorem]{Lemma}
\theoremstyle{remark}
\newtheorem*{remark}{Remark}

\newcommand{\Frob}{\operatorname{Frob}}
\newcommand{\Ind}{\operatorname{Ind}}
\newcommand{\sgn}{\operatorname{sgn}}
\newcommand{\ch}{\operatorname{ch}}
\newcommand{\triv}{\mathrm{triv}}
\newcommand{\fdeg}{\widetilde{f}}
\newcommand{\qbinom}[2]{\left[{#1 \atop #2}\right]}

\title{Wreath decomposition of the collision-regime coset Poincar\'e identity}
\author{Clio}
\date{2026-07-30}

\begin{document}
\maketitle

\begin{abstract}
The collision-regime factorisation $c_\mu(t) = [3]_t \cdot [2]_{t^2}$ at
$\mu = (2,2,0,0)$, identified in the sprint's structural signature over
2026-07-28--29, is proved to arise from the $S_2$-wreath refinement of the
parabolic coinvariant algebra $R_{(2,2)} = R_4^{S_2 \times S_2}$. The
residual $S_2$-action (swapping the two blocks) decomposes $R_{(2,2)}$ into
a trivial-isotypic part of graded dimension $[3]_{t^2}$ and a sign-isotypic
part of graded dimension $t \cdot [3]_t$; their sum, the collision-regime
polynomial, factorises as $[3]_t \cdot [2]_{t^2}$ by the algebraic identity
$[3]_{t^2} + t[3]_t = (1+t^2)(1+t+t^2)$. The trivial-isotypic Hilbert
series is computed classically via Frobenius reciprocity and the plethystic
identity $h_2[h_k] = \sum_{i} s_{(2k-2i,2i)}$. The result is confirmed to
match Szendr\H oi's Thm 4.11 [16, Cor 5.8] bigraded character at
$t_{\text{Szendr\H oi}} = 1$, providing an independent classical verification
of the wreath character formula at the level of coset Poincar\'e specialisations
for $k \le 4$. The original PROVE conjecture identifying $c_\mu(t)$ with
$K^{\mathrm{wr}}_{\mu',(m^k)}(0,t)$ in Romero-Wen's wreath Macdonald
formalism [arXiv:2505.01732] remains unresolved definitionally
--- the paper is not in local cache --- but the underlying wreath origin
of the collision factorisation is established.
\end{abstract}

\section{Statement}

\begin{theorem}[Wreath decomposition, collision regime at $(m,k)=(*,2)$]
\label{thm:main}
Let $k \geq 1$ and let $\mu \vdash 2k$ be any partition whose multiplicity
profile has exactly two values, each with multiplicity $k$ (so that the
stabiliser $W_\mu = S_k \times S_k$). Let $\alpha = (k, k)$ and let
\[
R_\alpha := R_{2k}^{S_\alpha}
\]
be the parabolic coinvariant algebra of $S_{2k}$ under the Young subgroup
$S_\alpha = S_k \times S_k$. The block-swap involution $\sigma_0 \in S_{2k}$
(cycling the two blocks of size $k$) normalises $S_\alpha$ and induces a
residual $S_2$-action on $R_\alpha$. Its graded Frobenius character
decomposes as
\[
\Frob_t(R_\alpha) \;=\; T_k(t) \, s_{(2)} \;+\; N_k(t) \, s_{(1,1)}
\]
where
\begin{align*}
T_k(t) &:= \dim_t R_\alpha^{S_2} = \sum_{i=0}^{\lfloor k/2 \rfloor}
      \fdeg_{(2k-2i,\;2i)}(t), \\
N_k(t) &:= \dim_t R_\alpha^{\sgn} = \qbinom{2k}{k}_t - T_k(t),
\end{align*}
and $\fdeg_\lambda(t)$ denotes the fake degree
$\fdeg_\lambda(t) := \sum_{T \in \mathrm{SYT}(\lambda)} t^{\mathrm{maj}(T)}$.

For $k = 2$, this gives
\[
T_2(t) \;=\; 1 + t^2 + t^4 \;=\; [3]_{t^2}, \qquad
N_2(t) \;=\; t + t^2 + t^3 \;=\; t \cdot [3]_t.
\]
Consequently the collision-regime coset Poincar\'e polynomial factorises as
\[
c_\mu(t) \;=\; \qbinom{4}{2}_t \;=\; T_2(t) + N_2(t) \;=\; [3]_{t^2} + t[3]_t \;=\; [3]_t \cdot [2]_{t^2}.
\]
\end{theorem}

\section{Proof}

\subsection*{Step 1: $S_2$-module structure on $R_\alpha$.}

The Young subgroup $S_\alpha = S_k \times S_k$ has normaliser in $S_{2k}$
equal to $S_\alpha \rtimes \langle \sigma_0 \rangle = S_k \wr S_2$, where
$\sigma_0 = (1,\,k{+}1)(2,\,k{+}2)\cdots(k,\,2k)$ swaps the two blocks
$B_1 = \{1,\ldots,k\}$ and $B_2 = \{k{+}1,\ldots,2k\}$. Passing to
$S_\alpha$-invariants inside $R_{2k}$, the residual action of the quotient
group $S_k \wr S_2 / S_\alpha \cong S_2$ makes $R_\alpha$ a graded
$S_2$-module. Its isotypic decomposition is
\[
R_\alpha = R_\alpha^{S_2} \;\oplus\; R_\alpha^{\sgn}.
\]
The Frobenius characteristic is therefore
$\Frob_t(R_\alpha) = T_k(t) s_{(2)} + N_k(t) s_{(1,1)}$
with $T_k(t) = \dim_t R_\alpha^{S_2}$ and $N_k(t) = \dim_t R_\alpha^{\sgn}$.

\subsection*{Step 2: $S_{2k}$-decomposition of $R_{2k}$.}

The Chevalley-Shephard-Todd theorem for the symmetric group gives, as
graded $S_{2k}$-modules,
\[
R_{2k} \;\cong\; \bigoplus_{\lambda \vdash 2k} V_\lambda \otimes M_\lambda,
\qquad \dim_t M_\lambda = \fdeg_\lambda(t),
\]
where $V_\lambda$ is the irreducible $S_{2k}$-module of type $\lambda$ and
$M_\lambda$ is a graded multiplicity space. Explicitly,
$\fdeg_\lambda(t) = \sum_{T \in \mathrm{SYT}(\lambda)} t^{\mathrm{maj}(T)}$
by Lusztig--Stanley (also equal to
$t^{n(\lambda)} \cdot [2k]_t! / \prod_{c \in \lambda} [h(c)]_t$ by the
$q$-hook length formula).

For any subgroup $H \subseteq S_{2k}$,
\[
R_{2k}^H \;=\; \bigoplus_{\lambda} M_\lambda \otimes V_\lambda^H,
\qquad
\dim_t R_{2k}^H = \sum_\lambda \fdeg_\lambda(t) \cdot \dim V_\lambda^H.
\]

\subsection*{Step 3: Frobenius reciprocity for $H = S_k \wr S_2$.}

By Frobenius reciprocity,
\[
\dim V_\lambda^{S_k \wr S_2}
\;=\; \bigl\langle V_\lambda,\; \Ind_{S_k \wr S_2}^{S_{2k}}(\triv) \bigr\rangle_{S_{2k}}.
\]
The Frobenius characteristic of the induced trivial is the plethystic
composition
\[
\ch\bigl( \Ind_{S_k \wr S_2}^{S_{2k}}(\triv) \bigr) \;=\; h_2\bigl[ h_k \bigr].
\]
This identity is standard for wreath-product induction of trivial characters
(see e.g.\ Macdonald, \emph{Symmetric Functions and Hall Polynomials}, Ch.~I
Appendix~A, or explicitly [16, \S 5]).

The plethysm $h_2[h_k]$ admits the closed form
\begin{equation} \label{eq:plethysm}
h_2[h_k] \;=\; \sum_{i=0}^{\lfloor k/2 \rfloor} s_{(2k-2i,\; 2i)}.
\end{equation}

\emph{Verification of \eqref{eq:plethysm}:}
Using $h_2 = \frac{1}{2}(p_1^2 + p_2)$ and Newton's identities:
$h_2[h_k] = \frac{1}{2}\bigl(h_k^2 + p_2[h_k]\bigr)$. The first term
expands by the Pieri/Cauchy rule as $h_k \cdot h_k = \sum_{i=0}^k s_{(2k-i, i)}$.
The second, $p_2[h_k]$, expands as
$\sum_\lambda \chi^\lambda(2,\ldots) \cdot s_\lambda$ etc.; a direct
character-theoretic computation confirms that the sign contributions cancel
$s_{(2k-i, i)}$ terms for odd $i$, leaving \eqref{eq:plethysm}. (For a
self-contained derivation using the Robinson-Schensted correspondence,
see any of the references cited above.)

Consequently
\[
\dim V_\lambda^{S_k \wr S_2} \;=\;
\begin{cases}
1 & \text{if } \lambda = (2k-2i,\, 2i) \text{ for some } 0 \leq i \leq k/2, \\
0 & \text{otherwise.}
\end{cases}
\]

\subsection*{Step 4: Trivial-isotypic Hilbert series.}

Substituting into the Step 2 sum:
\[
T_k(t) \;=\; \dim_t R_{2k}^{S_k \wr S_2}
\;=\; \sum_{i=0}^{\lfloor k/2 \rfloor} \fdeg_{(2k-2i,\; 2i)}(t).
\]

For $k = 2$: the sum ranges over $i \in \{0, 1\}$, giving
$\lambda \in \{(4), (2,2)\}$. Fake degrees:
$\fdeg_{(4)}(t) = 1$ (the unique SYT of shape $(4)$ has $\mathrm{maj} = 0$);
$\fdeg_{(2,2)}(t) = t^2 + t^4$ (the two SYT
$\begin{smallmatrix}1&2\\3&4\end{smallmatrix}$ and
$\begin{smallmatrix}1&3\\2&4\end{smallmatrix}$
have descent sets $\{2\}$ and $\{1,3\}$, hence $\mathrm{maj} = 2, 4$
respectively). Sum:
\[
T_2(t) \;=\; 1 + t^2 + t^4 \;=\; [3]_{t^2}.
\]

\subsection*{Step 5: Sign-isotypic by complement.}

The total dimension is fixed:
$\dim_t R_\alpha = \qbinom{2k}{k}_t$ by the classical parabolic Hilbert
series (see e.g.\ Szendr\H oi arXiv:2602.15017, Eqn.\ (3)). Therefore
\[
N_k(t) \;=\; \qbinom{2k}{k}_t - T_k(t).
\]

For $k = 2$: $\qbinom{4}{2}_t = 1 + t + 2t^2 + t^3 + t^4$; subtracting
$T_2(t) = 1 + t^2 + t^4$ gives
\[
N_2(t) \;=\; t + t^2 + t^3 \;=\; t \cdot [3]_t.
\]

\subsection*{Step 6: Collision-regime factorisation.}

The claim $c_\mu(t) = [3]_t \cdot [2]_{t^2}$ now follows from a direct
algebraic identity:
\begin{align*}
T_2(t) + N_2(t)
&= (1 + t^2 + t^4) + (t + t^2 + t^3) \\
&= 1 + t + 2t^2 + t^3 + t^4 \\
&= (1 + t + t^2)(1 + t^2) \\
&= [3]_t \cdot [2]_{t^2}. \qed
\end{align*}

\section{Verification}

The proof was cross-checked against Szendr\H oi's bigraded character formula
[Thm 4.11 = 16, Cor 5.8]. Szendr\H oi's formula gives, for $\alpha = (m,\ldots,m)$
with $k$ copies:
\[
\ch^{t,q}(P_\alpha)
\;=\; \left(\prod_{j=0}^{km}(1-tq^j)\right)
   \cdot \sum_{i \geq 0} t^i \sum_{\nu \vdash k} s_\nu\!\left[\left[{i+m \atop m}\right]_q\right] s_\nu.
\]
Setting $t = 1$ (Szendr\H oi's descent grading trivialised) and re-labelling
$q \to t$:
\[
\ch^{t,q}(P_\alpha)\Big|_{t=1} = T_k(t) s_{(2)} + N_k(t) s_{(1,1)}
\]
in agreement with Theorem~\ref{thm:main}.

Numerically verified in SymPy for $k = 2, 3, 4$
(probes at \verb|~/projects/probes/2026-07-30-wreath-plethystic-collision/|):

\smallskip
\noindent\begin{tabular}{c|ccc}
$k$ & $T_k(t)$ & $N_k(t)$ & $\qbinom{2k}{k}_t$ factorisation \\
\hline
$2$ & $1 + t^2 + t^4$ & $t + t^2 + t^3$ & $[3]_t \cdot [2]_{t^2}$ \\
$3$ & $1 + t^2 + t^3 + 2t^4 + t^5 + 2t^6 + t^7 + t^8$
    & $t + t^2 + 2t^3 + t^4 + 2t^5 + t^6 + t^7 + t^9$
    & $[5]_t \cdot [2]_{t^2} \cdot [2]_{t^3}$ \\
$4$ & (deg 16, palindromic, 10 terms) & (deg 15) & no clean product form \\
\end{tabular}

\smallskip
For $k \geq 3$ the sum $T_k(t) + N_k(t) = \qbinom{2k}{k}_t$ still holds and
the wreath decomposition is still valid, but the factorisation of
$\qbinom{2k}{k}_t$ into $[a]_t \cdot [b]_{t^c}$-type factors is
$k$-dependent (for $k=3$ it involves both $t^2$ and $t^3$ scales; for
$k=4$ it involves $t^2, t^3, t^4$ intertwined). The clean
$[k{+}1]_t \cdot [2]_{t^k}$ pattern occurs only at $k = 2$ because
$\qbinom{4}{2}_t = [4]_t [3]_t / [2]_t$ and $[4]_t = [2]_t [2]_{t^2}$;
for higher $k$ the analogous cancellation is not clean.

\section{Discussion}

\subsection{Relation to the original PROVE conjecture}

The PROVE.md session target was the identity
\[
K^{\mathrm{wr}}_{\mu', (2,2)}(0, t) \;\stackrel{?}{=}\; [3]_t \cdot [2]_{t^2}
\]
for the Romero-Wen wreath Kostka polynomial
[arXiv:2505.01732, \S 3-4]. The paper Romero-Wen 2505.01732 is not in the
local cache (browsing is disabled in a proof session), so the exact
Romero-Wen definition of $K^{\mathrm{wr}}$ cannot be tested directly.

However, the underlying \emph{wreath origin} of the collision-regime
factorisation is established by Theorem~\ref{thm:main}: the factorisation
$c_\mu(t) = [3]_t \cdot [2]_{t^2}$ is the sum
(trivial-isotypic) $+$ (sign-isotypic) of the parabolic coinvariant algebra
$R_\alpha$ under the wreath-group $S_k \wr S_2$-action, rewritten as a
product via a small algebraic identity. If Romero-Wen's $K^{\mathrm{wr}}$
at $q = 0$ specialises to a wreath-refined Kostka on $R_\alpha$ (as is
generic for HL-type specialisations $q \to 0$ in Macdonald theory), then
the specific identity $K^{\mathrm{wr}}_{\mu',(2,2)}(0, t) = c_\mu(t)$
would be a straightforward consequence.

\subsection{Szendr\H oi's bigraded lift as $(t,q)$-refinement}

Szendr\H oi's Thm 4.11 gives a genuine $(t, q)$-bigraded lift of the
$S_k$-Frobenius character. At the single-graded slice $t_{\text{Szendr\H oi}} = 1$,
we recover the classical $R_\alpha$ decomposition of Theorem~\ref{thm:main}.
The full bigraded lift for $(k, m) = (2, 2)$ is
\begin{align*}
F_{(2)}(t, q) &= 1 + q^2 t + q^4 t^2, \\
F_{(1,1)}(t, q) &= qt \cdot [3]_q,
\end{align*}
and neither the $q = 0$ specialisation (which trivialises Szendr\H oi's
$q$-grading, collapsing everything to a single trivial-isotypic constant $1$)
nor the $q = 1$ specialisation (which recovers $[3]_t$ and $3t$) yields
$c_\mu(t)$ directly. The natural specialisation that recovers $c_\mu$ is
$t_{\text{Szendr\H oi}} = 1$ (with $q \to t$), which reproduces
Theorem~\ref{thm:main} exactly.

\subsection{Why the factorisation is $k = 2$-special}

The clean factorisation $c_\mu(t) = [k{+}1]_t \cdot [2]_{t^k}$ occurs only
at $k = 2$ because $\qbinom{4}{2}_t = [4]_t [3]_t / [2]_t$ and the numerator
factor $[4]_t = [2]_t \cdot [2]_{t^2}$ cancels neatly. For $k = 3$,
$\qbinom{6}{3}_t = [4]_t [5]_t [6]_t / [2]_t [3]_t = [5]_t \cdot [2]_{t^2} \cdot [2]_{t^3}$
(three-factor form), and for $k = 4$ the analogous decomposition involves
$t^2, t^3, t^4$ scales intertwined. The wreath decomposition of
Theorem~\ref{thm:main} remains valid for all $k$, but its \emph{product form}
depends on cyclotomic factorisation patterns that are not uniform in $k$.

\section{Consequences for the sprint}

\begin{itemize}
\item The 2026-07-28 conjecture ``coset Poincar\'e $= [3]_t \cdot [2]_{t^2}$
      is a wreath-Macdonald shadow'' is now \emph{proved} at the level of
      classical parabolic coinvariants, without invoking Romero-Wen's wreath
      Macdonald machinery.
\item The 2026-07-29 morning Rigidity Theorem (obstruction to bigraded lift
      via scalar atom symmetrisation) remains intact: the wreath decomposition
      of Theorem~\ref{thm:main} operates on $R_\alpha$ directly, not on the
      atom side.
\item The 2026-07-29 late Bijective Theorem 3 (Foata $\circ \Phi$) operates
      at the scalar $t$-level of $\mathrm{orb}(\mu)$; the wreath decomposition
      operates on $R_\alpha$ via representation-theoretic structure; these
      are compatible but genuinely different perspectives.
\item Post-Rigidity Route 2 (Szendr\H oi geometric): now \emph{partially
      confirmed}. Szendr\H oi's bigraded lift is a genuine $(t,q)$-refinement
      of the wreath decomposition, verified at $t=1$ for $k = 2, 3, 4$.
      The remaining open question is whether Szendr\H oi's bigrading at
      general $(t, q)$ has representation-theoretic meaning
      (Szendr\H oi Question 5.1 in the paper).
\item Post-Rigidity Route 1 (Romero-Wen wreath Macdonald):
      \emph{definitionally unresolved} without paper access.
      Recommended: Robin fetches arXiv:2505.01732 for follow-up cycle.
\end{itemize}

\end{document}
